\section{Abel--Ruffini and finite Galois certificates}\label{sec:abel-ruffini}

The statement ``the quintic cannot be solved by radicals'' has two distinct
forms.  The generic theorem says that no single radical formula works for the
universal quintic.  A concrete theorem says that a specified irreducible
quintic has no radical expression for its roots.  Both can be organized
constructively, but neither says that every quintic is insoluble: reducible and
special solvable quintics are treated in Section~\ref{sec:solvable}.

\subsection{Radical towers yield solvable finite groups}

A finite radical certificate over a field \(K\) is a tower
\begin{equation}\label{eq:radical-tower}
 K=K_0\subset K_1\subset\cdots\subset K_m,
 \qquad K_{j+1}=K_j(\alpha_j),
 \qquad \alpha_j^{n_j}\in K_j,
\end{equation}
with the required roots of unity adjoined.  For object-theory purposes a finite
extension is represented by a basis, addition and multiplication tables,
embeddings of the preceding field, and reduction certificates for the defining
polynomials.  Every axiom is a bounded equality on table entries.

\begin{lemma}[Kummer step]\label{lem:kummer-step}
Suppose \(K\) contains all \(n\)-th roots of unity and
\(L=K(\alpha)\) with \(\alpha^n\in K\).  Every \(K\)-automorphism of the normal
closure sends \(\alpha\) to \(\zeta\alpha\) for an \(n\)-th root of unity
\(\zeta\).  The corresponding relative automorphism group embeds into the
abelian group of \(n\)-th roots of unity.
\end{lemma}

\begin{proof}
If \(\sigma\) fixes \(K\), then
\(\sigma(\alpha)^n=\sigma(\alpha^n)=\alpha^n\).  When \(\alpha\ne0\), the
quotient \(\sigma(\alpha)/\alpha\) is an \(n\)-th root of unity; the zero case
is trivial.  Multiplication of these quotients follows composition of
automorphisms.  In a coded finite extension, all divisions are implemented by
multiplicative-inverse tables with nonzero guards.
\end{proof}

\begin{theorem}[Radicals imply solvability]\label{thm:radicals-solvable}
The Galois group of the normal closure of a finite radical tower has a finite
normal series with abelian quotients.  Hence every polynomial solvable by
radicals has a solvable Galois group.
\end{theorem}

\begin{proof}
First adjoin the finitely many roots of unity required by the tower; these are
cyclotomic extensions with abelian Galois groups.  Normalize one radical step at
a time.  Lemma~\ref{lem:kummer-step} makes the kernel contributed by that step
abelian.  Induction on the tower length and the elementary fact that an
extension of a solvable group by a solvable group is solvable give the result.
For finite multiplication tables, normality, quotient multiplication, and the
derived series are all bounded checks.  The argument therefore translates into
\ipra{} for every coded tower.
\end{proof}

The converse---a polynomial with solvable Galois group is solvable by
radicals after adjoining roots of unity---uses the same series in reverse and
successive cyclic extensions.  In a fully explicit finite extension it is
again a certificate theorem.  We use both directions below for the standard
criterion ``solvable by radicals iff the finite Galois group is solvable.''

\subsection{The generic quintic has group \(S_5\)}

Let \(x_1,\ldots,x_5\) be algebraically independent indeterminates over
\(\Q\), let \(e_j\) be their elementary symmetric polynomials, and put
\[
 L=\Q(x_1,\ldots,x_5),
 \qquad K=\Q(e_1,\ldots,e_5).
\]
Every permutation of the \(x_i\) fixes \(K\), so \(S_5\) acts by
\(K\)-automorphisms of \(L\).  Conversely the roots of
\begin{equation}\label{eq:generic-quintic}
 Z^5-e_1Z^4+e_2Z^3-e_3Z^2+e_4Z-e_5
\end{equation}
are exactly the \(x_i\), so a \(K\)-automorphism permutes them.

The only subtle point is that the fixed field of the permutation action is
exactly \(K\).  This admits a finite algorithmic proof.

\begin{lemma}[Constructive symmetric rational functions]
\label{lem:symmetric-rational}
Every rational function in \(x_1,\ldots,x_5\) fixed by all permutations is a
rational function in \(e_1,\ldots,e_5\).  For bounded numerator and denominator
degrees, the conversion is primitive recursive.
\end{lemma}

\begin{proof}
Write the invariant as \(f=p/q\).  Replace the denominator by its permutation
norm
\[
 Q=\prod_{\sigma\in S_5}\sigma(q).
\]
Then \(Q\) is symmetric and \(P=fQ=p(Q/q)\) is a symmetric polynomial.
The fundamental theorem of symmetric polynomials has a terminating
lexicographic reduction: subtract from a symmetric polynomial a suitable
monomial in \(e_1,\ldots,e_5\) that cancels its leading monomial, and repeat.
The leading exponent strictly decreases in a finite well-order bounded by the
input degree.  Thus \(P=P(e_1,\ldots,e_5)\) and similarly for \(Q\).  The
reduction loop and its evaluation invariant are primitive recursive; the
termination measure is a bounded natural code of the exponent vector.
\end{proof}

\begin{theorem}[Generic Galois group]\label{thm:generic-S5}
The Galois group of the generic quintic~\eqref{eq:generic-quintic} over
\(K\) is \(S_5\).
\end{theorem}

\begin{proof}
The preceding discussion embeds the group in \(S_5\), while every permutation
of the indeterminates is a \(K\)-automorphism.  Lemma
\ref{lem:symmetric-rational} shows there is no larger fixed field hidden
between \(K\) and the invariant field, so the action is the full Galois action.
All coefficient manipulations at a requested degree are finite.  This is the
constructive core of the generic argument developed in ProVeIt's
polynomial-formulas/Abel--Ruffini packet.
\end{proof}

\subsection{Why \(S_5\) is not solvable}

Let the commutator convention be
\([g,h]=ghg^{-1}h^{-1}\).

\begin{lemma}\label{lem:S5-derived}
The commutator subgroup of \(S_5\) is \(A_5\).
\end{lemma}

\begin{proof}
The sign homomorphism has abelian target, so every commutator is even and
\([S_5,S_5]\subseteq A_5\).  On the other hand, the commutator of two
transpositions sharing one letter is a 3-cycle (up to orientation), and
3-cycles generate \(A_5\).  Hence \(A_5\subseteq[S_5,S_5]\).
\end{proof}

\begin{lemma}\label{lem:A5-perfect}
The group \(A_5\) is perfect: \([A_5,A_5]=A_5\).
\end{lemma}

\begin{proof}
Both
\[
 g=(2\,3)(4\,5),
 \qquad h=(1\,2)(4\,5)
\]
are in \(A_5\), and a direct five-entry permutation calculation gives
\[
 [g,h]=(1\,2\,3).
\]
The commutator subgroup is normal.  All 3-cycles are conjugate within
\(A_5\), and 3-cycles generate \(A_5\), so the commutator subgroup is all of
\(A_5\).
\end{proof}

Both lemmas are tiny finite certificates: the archive independently enumerates
the 120 permutations and obtains derived-series orders
\[
 120,60,60.
\]

\begin{theorem}[Constructive Abel--Ruffini theorem]
\label{thm:abel-ruffini}
There is no radical expression, uniform in the five generic coefficients, that
produces a root of the generic quintic.
\end{theorem}

\begin{proof}
By Theorem~\ref{thm:generic-S5}, the generic Galois group is \(S_5\).
Lemmas~\ref{lem:S5-derived} and~\ref{lem:A5-perfect} show that its derived
series never reaches the trivial group, so it is not solvable.  This contradicts
Theorem~\ref{thm:radicals-solvable} for any proposed radical tower.  Every group
calculation is finite, and the symmetric-polynomial reduction is a bounded
algorithm, so the argument has the required \ipra{} form.
\end{proof}

The theorem excludes a universal formula by ordinary radicals.  It does not
exclude Bring radicals, hypergeometric functions, numerical algorithms, or
radical formulas for special quintics.

\subsection{A finite group criterion for an irreducible quintic}

An irreducible quintic over \(\Q\) has a transitive Galois group
\(G\le S_5\).  The solvable transitive subgroups have a particularly small
form.

\begin{theorem}[Prime-degree solvable subgroup criterion]
\label{thm:AGL-criterion}
A solvable transitive subgroup of \(S_5\) is conjugate to a subgroup of
\begin{equation}\label{eq:AGL}
 \operatorname{AGL}(1,5)
 =\{x\mapsto ax+b:a\in\F_5^\times,\ b\in\F_5\},
\end{equation}
which has order \(20\).  Conversely every subgroup of
\(\operatorname{AGL}(1,5)\) is solvable.  Hence an irreducible quintic is
solvable by radicals exactly when its Galois group is conjugate to such a
subgroup.
\end{theorem}

\begin{proof}
Let \(G\) be solvable and transitive, and choose a nontrivial minimal normal
subgroup \(N\).  A standard finite induction on the derived series shows that
\(N\) is elementary abelian of prime-power order.  The orbits of a normal
subgroup form a block system.  Since the action has prime degree five, a
nontrivial normal subgroup must be transitive.  Therefore \(5\mid |N|\).
Inside \(S_5\), a 5-subgroup has order at most five, so \(N\cong C_5\) and acts
regularly.  Conjugation embeds a point stabilizer into
\(\Aut(C_5)\cong C_4\), giving
\(G\le C_5\rtimes C_4=\operatorname{AGL}(1,5)\) after relabeling the points.

Conversely the translation subgroup \(C_5\) is normal in
\(\operatorname{AGL}(1,5)\) and the quotient is \(C_4\), so the affine group
and all its subgroups are solvable.  The radical criterion follows from
Theorem~\ref{thm:radicals-solvable} and its converse.
\end{proof}

The executable verifier checks the theorem independently by enumerating all
156 subgroups of \(S_5\): there are 20 transitive subgroups as actual subsets,
18 solvable ones of orders \(5,10,20\), and two nonsolvable ones, \(A_5\) and
\(S_5\).  The affine group's derived-series orders are \(20,5,1\).

\subsection{A concrete nonsolvable quintic: \(x^5-x+1\)}

Let
\[
 f(X)=X^5-X+1.
\]
Its discriminant, from~\eqref{eq:trinomial-discriminant}, is
\[
 \Disc(f)=-256+3125=2869.
\]
Modulo two one has the exact square-free factorization
\begin{equation}\label{eq:f-mod2}
 f(X)=(X^2+X+1)(X^3+X^2+1)\pmod2.
\end{equation}
The quadratic and cubic factors have no root in \(\F_2\), hence are
irreducible.  The factor degrees are therefore \((2,3)\).

Modulo three, define \(r_0=X\) and
\(r_{j+1}=r_j^3\bmod f\).  A finite reduction gives
\begin{center}
\begin{tabular}{c|l}
\toprule
\(j\)&\(r_j\) in \(\F_3[X]/(f)\)\\
\midrule
1&\(X^3\)\\
2&\(-X^4+X-1\)\\
3&\(-X^4-X^2-1\)\\
4&\(-X^4-X^3+X^2+X-1\)\\
5&\(X\)\\
\bottomrule
\end{tabular}
\end{center}
Moreover the Bézout identity
\begin{equation}\label{eq:f-bezout3}
 f(X)+(-X^2-1)(X^3-X)=1\pmod3
\end{equation}
shows that \(f\) has no linear factor over \(\F_3\).  Since
\(f\mid X^{3^5}-X\) by the table, and five is prime, all irreducible factors
have degree five.  Thus \(f\) is irreducible modulo three and hence over
\(\Q\).

\begin{theorem}\label{thm:x5-x+1-S5}
The Galois group of \(X^5-X+1\) over \(\Q\) is \(S_5\).  In particular the
polynomial is not solvable by radicals.
\end{theorem}

\begin{proof}
The factorization modulo three supplies a Frobenius element of cycle type
\((5)\), while~\eqref{eq:f-mod2} supplies one of type \((2)(3)\).  Cubing the
latter leaves a transposition.  A subgroup of \(S_5\) containing a 5-cycle and
a transposition is all of \(S_5\): conjugating the transposition by powers of
the 5-cycle gives the edge transpositions of a connected graph on the five
letters, and edge transpositions of a connected graph generate the full
symmetric group.  Nonsolvability follows from
Lemmas~\ref{lem:S5-derived}--\ref{lem:A5-perfect}.
\end{proof}

The factor-degree/Frobenius lemma used here has a finite \ipra{} realization.
For a prime not dividing the discriminant, reduce the integral splitting order
modulo the prime, choose a prime ideal above it by bounded search in the finite
quotient, and use the residue Frobenius \(z\mapsto z^p\).  Its orbits on the
roots have lengths equal to the irreducible factor degrees.  Addition and
multiplication in every residue algebra are finite tables; the decomposition
and lifting checks are therefore certificate verifications.  This is a
substantial Lean module, but it needs no infinitary set existence.

\subsection{The real number \(\BR(3)\) is not a radical number}

Now let
\[
 g(X)=X^5+X+3.
\]
The discriminant is
\[
 \Disc(g)=256+3125\cdot3^4=253381.
\]
Modulo two, \(g\) has the same factorization~\eqref{eq:f-mod2}.  Modulo seven,
let \(s_0=X\) and \(s_{j+1}=s_j^7\bmod g\).  Exact reduction gives
\begin{center}
\begin{tabular}{c|l}
\toprule
\(j\)&\(s_j\) in \(\F_7[X]/(g)\)\\
\midrule
1&\(-X^3-3X^2\)\\
2&\(3X^4+2X^3+3X+1\)\\
3&\(2X^4-3X^3+3X^2+2X+3\)\\
4&\(2X^4+2X^3+X+3\)\\
5&\(X\)\\
\bottomrule
\end{tabular}
\end{center}
Writing \(h=s_1-X=-X^3-3X^2-X\), the explicit identity
\begin{equation}\label{eq:g-bezout7}
 -2g(X)+(-2X^2-X-2)h(X)=1\pmod7
\end{equation}
excludes linear factors.  As above, \(g\) is irreducible modulo seven.

\begin{theorem}\label{thm:BR3-not-radical}
The Galois group of \(X^5+X+3\) is \(S_5\).  Its unique real root
\(\BR(3)\) is not contained in any radical extension of \(\Q\).
\end{theorem}

\begin{proof}
Modulo seven the factor degree is five, giving a 5-cycle; modulo two the factor
degrees are two and three, giving a transposition after cubing.  Thus the group
is \(S_5\).  If the single algebraic number \(\BR(3)\) belonged to a radical
tower, the normal closure of that tower (after adjoining roots of unity) would
have solvable Galois group.  The splitting field of the minimal polynomial of
\(\BR(3)\) is a subextension of that normal closure, so its Galois group would
be a quotient of a solvable group and hence solvable, contradicting \(S_5\).
\end{proof}

This gives a concrete arithmetic interpretation of the source page's claim
that the Bring radical is not an ordinary radical function.  The generic
version follows already from Theorem~\ref{thm:abel-ruffini}; the specialization
at \(a=3\) shows that the obstruction is not merely formal.
